% Generated draft: requires mathematical and source review.
\documentclass[11pt]{article}
\usepackage[margin=25mm]{geometry}
\usepackage{fontspec}
\setmainfont{Latin Modern Roman}
\usepackage{amsmath,amssymb,mathtools,unicode-math}
\setmathfont{Latin Modern Math}
\usepackage{enumitem,needspace,xurl,hyperref}
\setlength{\parindent}{0pt}
\setlength{\parskip}{5pt}
\setlength{\emergencystretch}{3em}
\newcommand{\N}{\mathbb{N}}
\newcommand{\Z}{\mathbb{Z}}
\newcommand{\Q}{\mathbb{Q}}
\newcommand{\R}{\mathbb{R}}
\newcommand{\C}{\mathbb{C}}
\newcommand{\F}{\mathbb{F}}
\newcommand{\A}{\mathbb{A}}
\newcommand{\PP}{\mathbb{P}}
\newcommand{\E}{\mathbb{E}}
\newcommand{\eps}{\varepsilon}
\newcommand{\dd}{\,\mathrm{d}}
\begin{document}
\section*{UnsolvedMath ID: 30003970}
\textbf{Source problem code:} catalogue entry
\subsection{Definitions and mathematical statement}
For integers \(n\geq k\geq 2\), let \(K_n\) denote the complete graph on \(n\) vertices. Under the standard interpretation of the strong Ramsey game \(R(K_n,K_k)\), two players alternately claim one previously unclaimed edge of \(K_n\), with the first player moving first. A player wins immediately upon claiming every edge of a copy of \(K_k\). The game is drawn if all edges have been claimed without either player winning.

\textbf{Question.} Is it true that for every fixed integer \(k\geq 2\) there exist integers \(T(k)\geq 1\) and \(N(k)\geq k\) such that, for every integer \(n\geq N(k)\), the first player has a strategy in \(R(K_n,K_k)\) that guarantees a win on or before the first player's \(T(k)\)-th move, regardless of the second player's legal choices? In particular, the strategy must prevent the second player from winning earlier. The strategy may depend on \(n\) and \(k\), but the bound \(T(k)\) must be independent of \(n\).

\textbf{Interpretive qualifications.} The supplied statement names the strong Ramsey game without explicitly giving its rules or the range of \(k\); the formulation above uses the standard edge-claiming interpretation and the range \(k\geq 2\). The phrase ``as \(n\to\infty\)'' is interpreted as requiring the assertion for every sufficiently large \(n\). Moves are counted as moves of the first player; counting both players' turns instead gives the same boundedness question, since a win on the first player's \(t\)-th move occurs on turn \(2t-1\). These conventions require confirmation against the cited source.
\subsection{Short English statement}
For each fixed clique size \(k\), can the first player force a win in the strong edge-claiming game on \(K_n\) within a number of moves depending only on \(k\), for every sufficiently large \(n\)? The supplied record labels the problem open; this draft does not independently verify that status.
\subsection{Sources}
\begin{itemize}
\item Supplied local UnsolvedMath registry record: permanent ID \(30003970\), title \emph{Uniform Winning Times in Strong Ramsey Games}, source problem code \texttt{catalogue entry}, in the Oberwolfach Reports Open Problems collection.
\item Supplied primary-source citation: \emph{Mini-Workshop: Positional Games} (2019), Oberwolfach Reports, DOI \texttt{10.4171/owr/2018/44}, pp.\ 2715--2738. \url{https://doi.org/10.4171/owr/2018/44}. The background describes Oberwolfach Reports as documenting mathematical workshops and their open-problem discussions.
\item Supplied literature source: \emph{Strong Ramsey games}. \url{https://www.sciencedirect.com/science/article/pii/S0195669820300172}. No authors or further bibliographic details are supplied in the record.
\item Supplied literature assessment, checked on exactly \texttt{2026-08-21}: the record states that the cited primary literature still describes Beck's bounded-time complete-graph question as open. This is supplied dated curation derived from the OWR record, not an independently verified literature assessment. The record requests independent verification of time-sensitive status and cited claims.
\item Attribution and dates: the supplied assessment associates the question with Beck, while the explicit proposer field is empty. The record gives proposed year 2018 and the tag \texttt{owr-year:2018}, whereas its report citation is labeled 2019. These supplied date labels are retained without resolving their relationship.
\item Supplied editorial labels: problem status \texttt{open} and statement status \texttt{corrected\_verified}. The record reports manual semantic, typographic, and LaTeX review of the clean statement. These are source-record assertions and do not constitute independent mathematical or source verification of this draft.
\end{itemize}
\end{document}
